\chapter{Matriz de trabajo para agentes}

Este apéndice no está pensado como lectura principal del estudiante. Sirve para que el equipo de agentes mantenga el manual como una obra viva.

\section{Regla operativa}

Ningún material debe incorporarse al manual solo por su nombre de archivo. Antes debe convertirse en una ficha de conocimiento con:

\begin{itemize}
  \item problema;
  \item contexto;
  \item factores;
  \item niveles;
  \item unidad experimental;
  \item variables respuesta;
  \item variables controladas;
  \item diseño o método estadístico;
  \item evidencia disponible;
  \item limitaciones;
  \item uso didáctico;
  \item capítulo de integración.
\end{itemize}

\section{Ciclo de actualización}

\begin{enumerate}
  \item Ejecutar \texttt{python3 scripts/maestro.py investigar}.
  \item Revisar \texttt{memory/matriz\_investigacion\_material.md}.
  \item Seleccionar fichas útiles para cada capítulo.
  \item Redactar o modificar el archivo \LaTeX{} correspondiente.
  \item Compilar el manual.
  \item Revisar claridad, rigor y continuidad narrativa.
  \item Registrar correcciones en \texttt{memory/bitacora\_aprendizaje.md}.
\end{enumerate}

\section{Contrato de calidad}

Cada cambio debe mejorar al menos una dimensión:

\begin{itemize}
  \item claridad para el estudiante;
  \item integración de materiales reales;
  \item rigor metodológico;
  \item calidad de ejemplos;
  \item actividades de práctica;
  \item continuidad entre capítulos;
  \item posibilidad de compilar y editar.
\end{itemize}
